\section{\odag{} from scratch}
\label{sec:ontodag}

Now the other side. \odag{} \citep{ontodag} is a working system: a Python
package, a command-line tool \cmd{odag}, a web interface, an agent interface,
and a persistence layer that publishes to content-addressed storage. This
section develops it on the same six documents, with real transcripts.

\subsection{The shape of the thing}

\odag{} stores one graph. Nodes have names; edges mean ``is-under'', with the
more general node above. There is one distinguished node, the root
\cmd{*}, which is above everything. Edges may converge --- a node may have any
number of parents --- so the graph is a DAG rather than a tree, and by
Proposition~\ref{prop:dag-poset} it is a stored partial order.

The first design decision is a subtraction, and it is the one that most
distinguishes \odag{} from \fca{}:

\begin{keyidea}
\textbf{There is no distinction between objects and attributes.} Everything is
a node. \textsf{Flight} is a node; \texttt{outbound.pdf} is a node;
\textsf{Travel} is a node. A node may have children today and none tomorrow.
``Category'' and ``item'' are roles a node plays in a particular sentence, not
types in the data model.
\end{keyidea}

This is a deliberate contrast with both neighbours. The relational model
\citep{codd1970} fixes a schema per relation, which suits large collections of
homogeneous records rather than an open-ended universe of heterogeneous
things; description-logic ontology stacks \citep{gruber1993,owl2} draw a
firm line between classes and individuals and buy a great deal of
expressiveness with it. \odag{} declines both, and
Section~\ref{sec:one-sorted} weighs what that costs.

The consequence is that categories nest without any extra machinery.
\textsf{Flight} can be under \textsf{Travel} in the same way that
\texttt{outbound.pdf} is under \textsf{Flight}. In \fca{}, by contrast, an
attribute is an attribute forever: it lives in $M$, it cannot be an element of
$G$, and if you want a hierarchy \emph{of attributes} you must build it
outside the formalism or encode it by scaling. Section~\ref{sec:one-sorted}
develops this comparison properly.

\subsection{Building it: \texttt{put}}
\label{sec:put}

The write primitive is $\cmd{put}(x, \{p_1, \dots, p_k\})$: place node $x$
under each of the given parents, creating $x$ if needed. Here is the running
example, actually typed:

\begin{lstlisting}
$ odag put Travel '*'
$ odag put Japan '*'
$ odag put Booked '*'
$ odag put Refundable '*'
$ odag put Flight Travel
$ odag put Hotel Travel
$ odag put outbound.pdf  Flight Japan Booked Refundable
$ odag put inbound.pdf   Flight Japan Booked
$ odag put kyoto-inn.pdf Hotel  Japan Booked Refundable
$ odag put rail-pass.pdf        Japan Booked Refundable
$ odag put paris-hop.pdf Flight       Booked
$ odag put visa-note.md         Japan
\end{lstlisting}

\noindent (Silence means success: \cmd{odag} follows the Unix convention that
a tool says nothing when it works.)

For each requested edge $p \to x$, \cmd{put} performs four checks in order,
and the order matters:

\begin{enumerate}[leftmargin=1.7em]
\item \textbf{Already implied?} If $x$ is already reachable from $p$ by
  another path, the edge adds no information --- skip it silently.
\item \textbf{Would it create a cycle?} If $p$ is reachable from $x$, adding
  $p \to x$ would close a loop. Reject the \emph{whole operation} with an
  error, before modifying anything. This is what keeps the stored graph a
  partial order (antisymmetry $=$ acyclicity).
\item \textbf{Does it make an old edge redundant?} If some ancestor of $p$ has
  a direct edge to $x$, that edge is now implied by the new one --- delete it.
\item Add the edge.
\end{enumerate}

Steps 1 and 3 together maintain the transitive reduction incrementally. Step 3
is subtler than it looks: the redundancy can appear in either direction
(the new edge may bypass an existing edge from above \emph{or} from below),
and getting only one direction is a real bug that \odag{} shipped and later
fixed --- with the consequence, discussed in Section~\ref{sec:merge}, that
merge stopped being order-independent until it was corrected.

\begin{figure}[t]
\centering
\dotfig[0.72]{ontodag_travel}
\caption{The running example as an \odag{}, drawn from the live graph. Grey
nodes are the ones we would call categories, white ones the documents --- but
that is a distinction in the reader's head, not in the data. Arrows point from
special to general. Compare with the concept lattice of
Figure~\ref{fig:lattice}: same facts, different object.}
\label{fig:ontodag}
\end{figure}

\subsection{Querying: \texttt{get}}
\label{sec:get}

The entire query language is one operation.

\begin{definition}[The query primitive]
\[
  \cmd{get}(A_1, \dots, A_k) \;=\; \cone{A_1} \cap \cdots \cap \cone{A_k}
\]
--- everything that is below \emph{all} of the given categories.
\end{definition}

\begin{lstlisting}
$ odag get Japan Flight
inbound.pdf
outbound.pdf

$ odag get Japan Refundable
kyoto-inn.pdf
outbound.pdf
rail-pass.pdf

$ odag get Flight Hotel          # nothing is both
$ odag get Travel
Flight
Hotel
inbound.pdf
kyoto-inn.pdf
outbound.pdf
paris-hop.pdf
\end{lstlisting}

Three details are worth noticing in that transcript.

First, \cmd{get Travel} returns \textsf{Flight} and \textsf{Hotel} alongside
the documents. Cones do not stop at some notional item level, because there is
no item level. This is the one-sorted design showing through, and it is a
feature: ``show me everything travel-related'' should include the
sub-categories.

Second, \cmd{get Flight Hotel} returns nothing, and that is an answer, not an
error. An empty intersection is a perfectly good statement about the world.

Third, the empty query \cmd{get} with no arguments returns \emph{everything}:
the intersection of no cones is unconstrained, which makes the operation
total and makes \cmd{get}, \cmd{list} and \cmd{get '*'} one code path.

Two derived operations complete the query surface. \cmd{get\_any} takes
several such conjunctions and unions the results (disjunctive normal form,
computed entirely at query time, never stored). And \cmd{is\_below}$(x, y)$
answers the yes/no question directly, by walking \emph{upward} from $x$ ---
ancestor chains are short even in huge graphs, so this never touches a cone:

\begin{lstlisting}
$ odag below outbound.pdf Travel
true
$ echo $?
0
\end{lstlisting}

\paragraph{The planner.} A naive \cmd{get} would compute every cone and
intersect. \odag{} does better, in ways borrowed from database query
optimisers and all \emph{provably result-preserving}: it drops query terms
that are ancestors of other terms (they constrain nothing extra), orders the
remaining cones smallest-first using the exact \verb|descendant_count| stored
on every node, and then --- knowing at runtime how many candidates survive
each step --- chooses between \emph{walking} the next cone and \emph{probing}
upward from each surviving candidate. It stops the moment the running result
is empty. The maxim is: \emph{plan with what is known in advance, adapt on
what can only be learned by doing, and never let either change the answer.}

\begin{lstlisting}
$ odag count Japan
5
\end{lstlisting}

That \verb|descendant_count| is an exact, always-maintained
\verb|COUNT(*)| per category, and it is maintained by \emph{delta} rather than
recomputation --- appending a fresh node adds exactly one to each distinct
ancestor, contraction removes exactly one, and dropping a redundant edge
changes nothing by definition. It is also stored inside each record, hence
hashed, hence it must be an exact function of content and not of history.

\subsection{The golden rule, and canonical form}
\label{sec:canonical}

\begin{keyidea}
\textbf{Keep no link you can infer.} \odag{} stores, at all times, the
transitive reduction of the asserted relation. By
Theorem~\ref{thm:unique-reduction} that graph is \emph{unique}, so the stored
form depends only on \emph{what is known}, never on the order it was typed in
or on how many redundant links the user tried to add.
\end{keyidea}

This is easy to state and easy to underrate. It means the graph has a
\emph{canonical form}, and a canonical form can be fingerprinted.

\begin{measured}
Fifty random build orders of the same six facts, each replayed into a fresh
\odag{}: \textbf{1} distinct graph. Ten of those orders replayed through the
content-addressed persistence layer: \textbf{1} distinct root hash,
\begin{center}\ttfamily\footnotesize
5ff08fead6474d5e9cb1d0466c9a14a57f839f6e5d55813889c22bc44b0c700f
\end{center}
\end{measured}

\subsection{Editing: remove, move, and cone removal}

Three distinct operations, and conflating them loses data.

\cmd{remove}$(x)$ deletes a node and \emph{contracts} the graph over the hole:
$x$'s children are reattached to $x$'s parents, so nothing below $x$ becomes
unreachable. The result is again the unique minimal form.

\cmd{remove --cone}$(x)$ deletes $x$ \emph{and its contents}. In a
multi-parent graph ``its contents'' needs a definition, and the only sound one
is orphan collection: a member of the cone goes iff the root can no longer
reach it once $x$ is gone. So deleting a finished project removes the notes
that belonged only to it and spares the specification another project still
uses. Survivors are \emph{detached}, never reattached upward --- giving them
the deleted node's parents would assert something nobody said.

\cmd{move}$(x, \text{from } P, \text{to } Q)$ is the retracting twin of
\cmd{put}: assert the new parents, retract the old. Everything below $x$
travels with it for free, because membership is reachability rather than a
stored subtree. A move can leave a shared child under both the old and the new
category, which is true rather than broken --- subsumption inherits, exclusive
status cannot --- so the tool reports those items instead of silently
deciding.

\subsection{Merging: the CRDT property}
\label{sec:merge}

\cmd{merge}$(D_1, D_2)$ folds one graph into another: add every missing node,
replay every edge through the same \cmd{put} machinery so that redundancies
created by the union are pruned, and the result is again canonical.

\begin{theorem}[Merge algebra; tested as invariant I7 in \citealp{ontodag}]
\label{thm:merge}
\cmd{merge} is commutative and idempotent: merging $A$ into $B$ gives the same
graph as merging $B$ into $A$, and merging the same thing twice changes
nothing.
\end{theorem}

Those two properties, plus the canonical form, are the definition of a
state-based conflict-free replicated data type \citep{shapiro2011}. Their
practical meaning:

\begin{keyidea}
Any group of people can exchange changes \textbf{in any order, repeatedly,
with no coordination}, and converge on the byte-identical graph. No server, no
consensus protocol, no locking. Two writers who have folded in each other's
work agree, and they can prove they agree by comparing 32 bytes.
\end{keyidea}

Why does it work? Because the canonical state is defined as \emph{the
transitive reduction of the union of asserted edges}, and union is
commutative and idempotent while reduction is a function. The correctness of
this argument depends on reduction being \emph{complete} --- which is why the
one-directional-redundancy bug mentioned in Section~\ref{sec:put} was not
cosmetic: with reduction incomplete, two writers could land on different roots
for the same knowledge, and Theorem~\ref{thm:merge} was false.

The price is stated openly: merge is union, so it is \emph{grow-only}. A
removal can be resurrected by a peer who still holds the old edge. Deletion
that must propagate needs either tombstones or a social protocol; \odag{}
declines both and documents the wall.

\subsection{Typed values: the order that is computed}
\label{sec:dimensions}

Everything so far is asserted edges. But no finite set of edges can say that a
parcel weighing $4\,\mathrm{kg}$ falls under ``at most $5\,\mathrm{kg}$'',
because between any two weights there is another. So for \emph{parametric}
terms, the order is \emph{computed from the name} instead of stored.

A term like \verb|weight(4kg)| denotes a set of values; \verb|weight(..5kg)|
denotes the set of weights at most five kilograms; and one term is below
another exactly when its denoted set is contained in the other's --- the same
``is a special case of'' meaning every stored edge has. Values are exact
rationals of the SI anchor unit, so every comparison is exact arithmetic and
every replica computes the identical order.

\begin{lstlisting}
$ odag prelude                       # declare the standard dimensions
$ odag put cello 'weight(4kg)'
$ odag put crate 'weight(25kg)'
$ odag put flute 'weight(400g)'

$ odag get 'weight(..5kg)'           # everything under five kilos
cello
flute
weight(2/5kg)
weight(4kg)

$ odag canon 'weight(400g)'          # what is actually stored
weight(2/5kg)

$ odag --render get 'weight(..5kg)'  # the same answer, read back nicely
cello
flute
weight(400g)
weight(4kg)
\end{lstlisting}

Three things in that transcript matter for the comparison with \fca{}.

First, \verb|weight(..5kg)| is a \emph{virtual} query term: no such node
exists, and none is created. The query is answered by filtering the values
that do exist. Constraints cost no writes.

Second, \verb|400g| is stored as \verb|weight(2/5kg)| --- canonicalised by
\emph{denotation}, not by spelling. Two people who file the same weight in
different units produce the identical stored form and hence the identical
hash. This is the semantic canonicalisation of Section~\ref{sec:canonical}
extended to values.

Third, the round trip \verb|--render| shows the readable spelling without ever
storing it. The user's phrasing is deliberately \emph{not} kept: if it were,
two people who know the same fact would disagree on the root.

\begin{remark}
Section~\ref{sec:learn-patterns} shows that this construction is an instance
of a known \fca{} generalisation --- interval pattern structures
\citep{ganter2001,kaytoue2011} --- arrived at independently, and that the
theory has things to say about it.
\end{remark}

\subsection{Persistence: from a graph to 32 bytes}
\label{sec:persistence}

\odag{}'s persistence goes through \texttt{recordstore} \citep{recordstore}, a
small library providing a versioned key$\to$record store over
\emph{content-addressed} storage --- storage where a blob is retrieved by the
hash of its own bytes and nothing is ever overwritten.

One record per node, keyed by name, holding its parents, children, count and
optional payload. The index over those records is a canonically encoded
radix trie, built so that equal contents always produce an equal index. So
\verb|commit()| returns one hash, the \emph{root}, with three properties:

\begin{itemize}[leftmargin=1.4em]
\item the root names the entire dataset at that moment --- share the root,
  share the dataset; keep an old root, keep a snapshot forever;
\item equal content gives an equal root, always, so comparing two versions
  starts by comparing two strings;
\item consecutive versions \emph{share structure}: a small change writes only
  the changed records plus a thin path of index nodes.
\end{itemize}

Pointing the same machinery at a peer-to-peer content-addressed network
(Ethereum Swarm; \citealp{tron2020}) publishes the ontology into that network,
with the latest root announced through a signed feed that others can follow.
Section~\ref{sec:crypto} is about what this buys.

\subsection{Residency: eager, lazy, sparse}

The same graph semantics come in three residencies, which matters for the
thin-client story in Section~\ref{sec:agents}. \texttt{EagerOntoDAG} loads
every record into memory. \texttt{LazyOntoDAG} \emph{reads} a published graph
by fetching records only as a query walks them, so querying a published store
costs the query rather than the store. \texttt{SparseOntoDAG} \emph{writes}
the same way: an edit touches only the records the reduction rules must
check, and commits the byte-identical root a fully loaded editor would have
produced. Published \emph{cone summaries} --- one derived record stating a
broad category's membership, in a separate store with its own root --- turn
broad queries from hundreds of fetches into a handful.

\subsection{The guarantees}
\label{sec:guarantees}

\odag{} publishes a contract of what a higher layer may assume. Six clauses,
each backed by a conformance test:

\begin{center}\small
\begin{tabularx}{\textwidth}{@{}lX@{}}
\toprule
\textbf{G1 Canonical root} & Equal knowledge yields an equal root. The stored
form is a \emph{semantic} canonical form: build history, insertion order and
the spelling of equal denotations do not affect it.\\
\textbf{G2 Monotone merge} & Merge is union then re-reduction. True
\cmd{is\_below} answers stay true; query results only grow. (Exception: a
\cmd{remove} loses to a concurrent re-add.)\\
\textbf{G3 Determinism} & Same root and same interpretation context give the
same answers on any replica and any residency.\\
\textbf{G4 Fail-closed \cmd{is\_below}} & True only when the graph plus exact
arithmetic \emph{witness} it. False means ``not derivable'', never ``unknown
but plausible''.\\
\textbf{G5 Convergence} & Writers who fold in each other's published roots
reach byte-identical roots regardless of gossip order.\\
\textbf{G6 Overlap completeness} & \cmd{get\_overlapping} returns a
recall-complete candidate set for a parametric term: complete for
possibility, silent on satisfaction.\\
\bottomrule
\end{tabularx}
\end{center}

Notice how many of these are \emph{negative} guarantees --- promises about
what will \emph{not} happen (no history dependence, no shrinking answers, no
plausible-but-underived \emph{true}). Section~\ref{sec:agents} argues that
negative guarantees are exactly what a machine consumer needs, and exactly
what a language model cannot give about its own knowledge.

\subsection{Summary of Section~\ref{sec:ontodag}}

\begin{keyidea}
\odag{} is \textbf{assertion kept minimal}. One graph, one root, no
object/attribute distinction, one query primitive (intersection of cones),
stored always in its unique transitive reduction. That uniqueness makes the
graph canonical, canonicity makes it hashable, hashability makes it
shareable and verifiable, and the union-plus-reduction definition of merge
makes concurrent writers converge without coordination.

It answers the question \emph{``what has been said, and what follows from it
right now?''} --- cheaply, incrementally, and in a form two strangers can
compare in 32 bytes.
\end{keyidea}
